

\begin{exercise}[Basic application of continuous mapping theorem \Coffeecup]
    Let \((X_n)_{n\in \nat}\) and \(X\) be real-valued random variables
    with \(X_n \overset{d}\to X\). Show that
    \begin{enumerate}[label=(\alph*), noitemsep]
        \item\label{it: (basic appl cont map) square}
        \(X_n^2 \overset{d}\to X^2\),

        \item\label{it: (basic appl cont map) sign}
        \(\sgn(X_n) \overset{d}\to \sgn(X)\) if \(\prob{X=0} = 0\).

        Construct a counterexample for the case \(\prob{X=0} > 0\).
    \end{enumerate}
\end{exercise}

\begin{solution}
\begin{enumerate}[wide=0pt]
    \item[\ref{it: (basic appl cont map) square}]
    This is a direct application of the continuous mapping theorem with the function \(f(x) = x^2\) that is continuous everywhere.            

    \item[\ref{it: (basic appl cont map) sign}]
    For the first part we again apply the continuous mapping theorem with the function \(f(x) = \sgn(x)\) that is continuous at all points except zero. Since \(\prob{X=0} = 0\) we have that \(f\) is almost surely continuous in \(X\) and thus the continuous mapping theorem applies.

    For the counterexample we can take any deterministic sequence \(x_n\to 0\)
    such that \(\sgn(x_n)\) does not converge to \(\sgn(0)\). This does depend a bit on the definition
    of \(\sgn(0)\) but exists in all cases. Then \(X_n = x_n\) has the same problem.
    Specifically, for the Lipschitz function \(g(x) = \min\set{\abs{x},
    1}\) we have
    \[
        \E{g(X_n)} = x_n \not\to 0 = \E{g(0)}
    \]
    and thus \(X_n \overset{d}{\not\to} 0\) by \ref{it: b.lip}.
    \qedhere
\end{enumerate}
\end{solution}

